\section{Introduction}

Modern recognition models are trained on large-scale datasets of natural images. Such images capture objects in their typical surroundings, \eg, a cow on grass, a boat in water, or a stove in a kitchen. These models can exploit co-occurrence cues that go unnoticed on standard benchmarks, since training and test sets draw from the same patterns. Failures surface only when objects appear in atypical contexts, \eg, a cow on a beach~\cite{beery2018recognition}. While context is itself informative for recognition~\cite{torralba2003contextual,oliva2007role}, models can come to use it as a shortcut around the object's visual features. This \emph{over-reliance} on co-occurring context has been studied through context-based interventions~\cite{shetty2019not,singh2020don,chakraborty2024visual,sarridis2025mavias}, controlled benchmarks~\cite{liu2015deep,sagawa2020distributionally,li2023whac}, and dedicated out-of-context detection~\cite{acharya2022detecting,bomatter2021pigs}. All these studies converge on the same conclusion: classifiers exploit co-occurrence shortcuts and fail when they are absent.

\looseness=-1It is unclear whether these findings still hold across the modern vision stack. Recent architectures and training paradigms such as end-to-end detection transformers~\cite{carion2020end,zong2023detrs,zhang2022dino}, open-vocabulary detectors that ground predictions in language~\cite{liu2024grounding,minderer2022simple,minderer2023scaling}, self-supervised vision foundation models~\cite{caron2021emerging,oquab2024dinov,simeoni2025dinov3}, vision-language models pretrained on billions of image-text pairs~\cite{radford2021learning,cherti2023reproducible,zhai2023sigmoid,tschannen2025siglip}, and large multimodal models that fuse vision and language at scale~\cite{li2024llava} are expected to alleviate context reliance through enhancements such as deformable attention, language grounding, as well as web-scale data diversity~\cite{xia2022vision,liu2024grounding,radford2021learning}. Recent work has examined whether such shortcuts persist for specific subsets of the stack: segmentation and detection through explainable artificial intelligence (XAI) methods~\cite{dreyer2023revealing}, and nuisance-based robustness benchmarks~\cite{zhao2022ood,dunkel2025cns}. However, a systematic diagnostic across diverse model families that evaluates on object pairs never co-occurring in natural data remains underexplored.

\begin{figure}[t]
\centering
\includegraphics[width=\linewidth]{static/images/teaser_tex.pdf}

\vspace{-3pt}
\caption{\textbf{Typicality bias across the modern vision stack.}
\emph{Left:} Three \emph{anchor classes} (target object categories we examine): \textit{\textcolor{anchorcolor}{skateboard}}, \textit{\textcolor{anchorcolor}{dog}}, and \textit{\textcolor{anchorcolor}{cow}}, shown in their {\color{typcolor}\textbf{typical}} company (top row: with a \textit{\textcolor{typcolor}{person}}, with a \textit{\textcolor{typcolor}{frisbee}}, with a \textit{\textcolor{typcolor}{bird}}) and in {\color{atypcolor}\textbf{atypical}} surroundings (bottom row: with a \textit{\textcolor{atypcolor}{cat}}, with a \textit{\textcolor{atypcolor}{spoon}}, with a \textit{\textcolor{atypcolor}{clock}}). \emph{Right:} Each model's score on the anchor classes in COCO \texttt{val} images where they appear with typical vs.\ atypical partners (NPMI quintile splits, \S\ref{sec:diagnostic}). The performance gap $\textcolor{atypcolor}{\Delta} =$ (\textcolor{black}{Typical} -- \textcolor{black}{Atypical}) persists across all five model families and three tasks we examine, ranging from $\textcolor{atypcolor}{\Delta = 5.4}$~p.p.\@ (Co-DINO Swin-L) to $\textcolor{atypcolor}{\Delta = 11.0}$~p.p.\@ (OpenCLIP).}
\label{fig:teaser}
\vspace{-12pt}
\end{figure}


In this work, we revisit this phenomenon systematically and at scale across modern vision models. We refer to it as \textit{\textbf{typicality bias}}: the tendency of recognition models to base predictions on familiar co-occurring context rather than the target object's own visual evidence. As shown in Fig.~\ref{fig:teaser}, a model exhibiting typicality bias performs well when objects appear in typical contexts and degrades when they appear alone or in atypical contexts.

To diagnose typicality bias systematically, we introduce a framework based on normalized pointwise mutual information (NPMI) over object pairs. For each anchor class, we use NPMI to construct paired evaluation splits in which the anchor (object category under examination) appears alongside either typical or atypical partners. To extend this natural-data evaluation with co-occurrence regimes that natural data do not cover (see Fig.~\ref{fig:npmi_matrix}), we further release \textbf{STRANGE-Bench} (\textbf{S}ynthetic \textbf{T}ypicality-\textbf{RAN}ked \textbf{G}rounded \textbf{E}valuation \textbf{Bench}mark), a photorealistic, layout-controlled synthetic benchmark spanning the full spectrum from typical to never-co-occurring object pairs.

~\looseness=-1A complementary question is how to mitigate typicality bias. Data augmentation is a natural lever for this, since it directly shapes the training distribution that produces the bias. In practice, existing augmentations are designed to add visual diversity, not to break co-occurrence shortcuts: both non-generative~\cite{bochkovskiy2020yolov4,zhang2017mixup,ghiasi2021simple} and generative~\cite{zhao2023x,kupyn2024dataset,xie2025mosaicfusion} methods sample from the natural training distribution, preserving rather than perturbing it. We present the first systematic study of how popular data augmentations interact with typicality bias across representative model families, and propose a simple \emph{context-aware} sampling rule that reduces the bias while preserving overall accuracy.

Our findings reveal three patterns. \textit{First,} typicality bias persists across diverse model families and tasks, with gaps of $\mathbf{5}$--$\mathbf{11}$~p.p.\@ between typical and atypical splits. \textit{Second}, the bias does not consistently shrink with scale: more pretraining data does not narrow the gap, and larger models narrow it without closing it, indicating that scale alone is not a solution. \textit{Third}, most non-paste augmentations leave the gap largely unchanged or narrow it by lowering both sides, while paste-based ones narrow it on STRANGE-Bench. AAS lifts atypical scenes further without hurting typical ones, and reduces the gap of standard training by up to $\mathbf{1.6}$~p.p.\@ on natural-image COCO~\cite{mscoco14} splits and $\mathbf{23.0}$~p.p.\@ on synthetic STRANGE-Bench, without sacrificing overall performance.

Our main \textbf{contributions} are summarized as follows:

\begin{enumerate}[leftmargin=*,nosep]
\item \textbf{Typicality bias persists across the modern vision stack.} We provide a unified diagnostic of typicality bias spanning five model families (specialist detection, open-vocabulary detection, self-supervised representations, vision-language models, and multimodal LLMs) and three recognition tasks. Our key finding is that typicality bias persists across diverse tasks and scales.

\item \textbf{A photorealistic synthetic benchmark for tiered typicality.} We present \textbf{STRANGE-Bench}, a layout-controlled synthetic benchmark spanning the full spectrum from typical to never-co-occurring object pairs, enabling controlled evaluation in low-NPMI regimes that natural data cannot cover; blinded user studies confirm that its gaps are not rendering artifacts.

\item \textbf{Anchor-Atypicality Sampling for mitigation.} We present the first systematic study of how popular data augmentations interact with typicality bias across representative model families, and propose Anchor-Atypicality Sampling (AAS), a simple sampling rule for paste-based augmentation that reduces typicality bias while preserving overall performance.
\end{enumerate}
